\documentclass{article}

\usepackage{amsmath,amssymb}
\usepackage{xurl}
\usepackage[hidelinks]{hyperref}
\setlength{\emergencystretch}{2em}

% Shared commands for notes in docs/paper-gaps/.

\DeclareUnicodeCharacter{2080}{\ifmmode{}_{0}\else\textsubscript{0}\fi}
\DeclareUnicodeCharacter{2081}{\ifmmode{}_{1}\else\textsubscript{1}\fi}
\DeclareUnicodeCharacter{2082}{\ifmmode{}_{2}\else\textsubscript{2}\fi}
\DeclareUnicodeCharacter{2099}{\ifmmode{}_{n}\else\textsubscript{n}\fi}

\newcommand{\Lean}{\textsc{Lean}}
\ExplSyntaxOn
\NewDocumentCommand{\leanid}{m}{
  \ifmmode
    \textsf{\detokenize{#1}}
  \else
    \group_begin:
    \urlstyle{sf}
    \tl_set:Nn \l_tmpa_tl {#1}
    \tl_replace_all:Nnn \l_tmpa_tl {\_} {_}
    \exp_args:NV \path \l_tmpa_tl
    \group_end:
  \fi
}
\ExplSyntaxOff
\providecommand{\tr}{\operatorname{tr}}
\newcommand{\tnleanIssueBase}{https://github.com/LionSR/TNLean/issues/}
\newcommand{\tnleanPrBase}{https://github.com/LionSR/TNLean/pull/}
\newcommand{\tnleanDocBase}{https://sirui-lu.com/TNLean/docs/}
\newcommand{\ghissue}[1]{\href{\tnleanIssueBase#1}{GitHub issue \##1}}
\newcommand{\ghpr}[1]{\href{\tnleanPrBase#1}{PR \##1}}
\newcommand{\doclink}[1]{\href{\tnleanDocBase#1.html}{\path{#1}}}

% Dirac notation and the identity operator, as in blueprint/src/macros/common.tex.
% \providecommand keeps note-local definitions authoritative.
\usepackage{dsfont}
\providecommand{\ket}[1]{|#1\rangle}
\providecommand{\bra}[1]{\langle #1|}
\providecommand{\braket}[2]{\langle #1 | #2 \rangle}
\providecommand{\ketbra}[2]{|#1\rangle\!\langle #2|}
\providecommand{\Id}{\mathds{1}}
\providecommand{\one}{\mathds{1}}
\providecommand{\C}{\mathbb{C}}
\providecommand{\MN}[1]{M_{#1}(\C)}
\providecommand{\MD}{M_D(\C)}

% External referencing for paper sources.  The build helper writes sanitized
% auxiliary files under build/paper-aux/ and excludes duplicated source labels.
\usepackage{xr}

% Import a generated paper auxiliary file with a caller-supplied label prefix.
% The candidate paths support builds from docs/paper-gaps/, the repository root,
% and build/paper-aux/ itself.
\newcommand{\importpaperaux}[2]{%
  \IfFileExists{../../build/paper-aux/#2.aux}{%
    \externaldocument[#1]{../../build/paper-aux/#2}%
  }{%
    \IfFileExists{build/paper-aux/#2.aux}{%
      \externaldocument[#1]{build/paper-aux/#2}%
    }{%
      \IfFileExists{#2.aux}{%
        \externaldocument[#1]{#2}%
      }{}%
    }%
  }%
}

% General external reference macros
\newcommand{\paperref}[2]{\ref{#1-#2}}
\newcommand{\papereqref}[2]{(\ref{#1-#2})}

% CPSV16 source (1606.00608 / MPDO-22-12-17-2.tex) external document and shortcuts
\importpaperaux{cpsv16-}{1606.00608/MPDO-22-12-17-2-tnlean}
\newcommand{\cpsvref}[1]{\paperref{cpsv16}{#1}}
\newcommand{\cpsveqref}[1]{\papereqref{cpsv16}{#1}}

% Verdict marker: \gapnote{<kind>}{<status>}. Kind is the policy
% classification (clarification, local-correction, scope-restriction,
% unfaithful, false-source, open-gap); status is open, wip (elimination
% underway), resolved, or historical. Machine-read by the site index;
% prints a verdict line.
\newcommand{\gapnote}[2]{\par\noindent\textbf{Verdict:} #1 (#2).\par}


\title{Sharp Finite-Length Block Separation for Canonical-Form Tensors}
\author{}
\date{2026-07-15}

\begin{document}
\maketitle
\gapnote{scope-restriction}{resolved}

\section{The assertion}

This note concerns the sharp block-injectivity proposition of Cirac,
P\'erez-Garc\'ia, Schuch, and Verstraete
\cite[lines~340--345]{Cirac2016MPDO_arXiv}.\footnote{For traceability,
    this note corresponds to
    \href{https://github.com/LionSR/TNLean/issues/1043}{issue \#1043}.
    Related formal work is recorded in
    \href{https://github.com/LionSR/TNLean/issues/587}{issue \#587} and
    \href{https://github.com/LionSR/TNLean/issues/822}{issue \#822};
    earlier proof attempts are
    \href{https://github.com/LionSR/TNLean/pull/682}{PR \#682} and
    \href{https://github.com/LionSR/TNLean/pull/813}{PR \#813}.}

Let $A$ be a canonical-form tensor with total bond dimension $D$. Let
$A_1,\ldots,A_g$ be a basis of normal tensors (BNT) for $A$, where $A_j$ has
bond dimension $D_j$. The source characterizes this basis by minimality: two
distinct representatives cannot be related by an invertible similarity and a
phase \cite{Cirac2016MPDO_arXiv}. For a physical word $w$, write $A_j^w$ for
its evaluation on the $j$th representative. Define the product algebra by
\begin{align}
    \mathcal A_{\mathrm{prod}}
        &= \bigoplus_{j=1}^{g} M_{D_j}(\C).
    \label{eq:cpgsv17_bicf_block_separation_product_algebra}
\end{align}
The tensor is in block-injective canonical form when its simultaneous word
evaluations span $\mathcal A_{\mathrm{prod}}$. Equivalently, for every tuple
$X=(X_j)_{j=1}^g\in\mathcal A_{\mathrm{prod}}$, a linear combination of
physical words evaluates to $X_j$ on every block $j$.

The source first applies the quantum Wielandt theorem to make each normal
representative injective after at most $D^4$ sites
\cite{Sanz2010Quantum,Cirac2016MPDO_arXiv}. It then uses the cross-block
argument of P\'erez-Garc\'ia, Verstraete, Wolf, and Cirac
\cite{PerezGarcia2007MPSReps}: if every normal block is injective at length
$L$, then the tensor is block injective after $3(L+1)(D-1)$ sites. The sharp
proposition therefore states that blocking at most $3D^5$ spins places every
canonical-form tensor in block-injective canonical form
\cite[lines~340--345]{Cirac2016MPDO_arXiv}. The relevant definition appears at
lines~318--322 of the same source.

\section{The older per-copy formulation}

The older auxiliary hypotheses impose finite-length separation on every
flattened copy. Let $\mathcal W_N$ denote the set of physical words of length
$N$. The trace-separation hypothesis states that there is an $N$ such that,
for matrices $\Delta_j\in M_{D_j}(\C)$,
\begin{align}
    \left(\sum_{j=1}^{g}\tr\!\left(\Delta_j A_j^w\right)=0
        \text{ for every }w\in\mathcal W_N\right)
    &\Longrightarrow
    \left(\Delta_j=0\text{ for every }j\right).
    \label{eq:cpgsv17_bicf_block_separation_trace_separation}
\end{align}
The trace pairing on $\mathcal A_{\mathrm{prod}}$ is nondegenerate, so
Eq.~\ref{eq:cpgsv17_bicf_block_separation_trace_separation} is dual to the
finite-length spanning condition
\begin{align}
    \operatorname{span}
    \left\{\left(A_1^w,\ldots,A_g^w\right)
        : w\in\mathcal W_N\right\}
    &= \mathcal A_{\mathrm{prod}}.
    \label{eq:cpgsv17_bicf_block_separation_word_tuple_span}
\end{align}

Three exact hypotheses used in the formal development imply
Eq.~\ref{eq:cpgsv17_bicf_block_separation_trace_separation}. The first is the
product-algebra span in
Eq.~\ref{eq:cpgsv17_bicf_block_separation_word_tuple_span}. The second requires
linear independence, at one common length, of the scalar functions
$w\mapsto(A_j^w)_{\alpha\beta}$ indexed by the block $j$ and the matrix entry
$(\alpha,\beta)$. This entrywise condition implies the full product-algebra
span. The third is a selector-word version of sharp block injectivity: each
block is injective after a common blocking length, and a further finite length
provides word polynomials that equal the identity on one selected block and
vanish on every other block. Concatenating an injective prefix with a selector
suffix gives Eq.~\ref{eq:cpgsv17_bicf_block_separation_word_tuple_span}.

These criteria yield the older per-copy trace-separation conclusion in the
formal development.\footnote{The criteria are
    \leanid{MPSTensor.WordTupleSpanTop}, \leanid{MPSTensor.wordEntryFamily}, and
    \leanid{MPSTensor.PropBlockInjective}. The trace-separation assumption is
    \leanid{MPOTensor.HorizontalCFData.biCF}. The constructors from these hypotheses to
    \leanid{MPOTensor.HorizontalCFData} structures are in the modules under
    \path{TNLean/MPS/MPDO/BiCFDerivation/}. Relevant implications include
    \leanid{MPSTensor.hasBiCF_of_wordTupleSpanTop},
    \leanid{MPSTensor.hasBiCF_of_wordEntryFamily_linearIndependent},
    \leanid{MPSTensor.hasBiCF_of_propBlockInjective}, and corresponding
    \leanid{MPOTensor.HorizontalCFData} constructors such as
    \leanid{MPOTensor.horizontalCFData_of_propBlockInjective}. The blueprint records
    these restricted criteria in its finite-witness remark.}
They remain valid restricted results. The sharp proposition is now proved
separately for a representative-indexed BNT sector decomposition, which groups
repeated copies before applying the separation argument.

\section{Why the fields are stated independently}

Blockwise injectivity, left-canonicality, and nonzero copy weights do not by
themselves imply the finite-length trace-separation witness. Consider two
one-dimensional scalar blocks over a one-letter alphabet. Let both block
tensors equal $1$, with respective weights $1$ and $2$. At every length, the
simultaneous word tuples span only
\begin{align}
    \operatorname{span}\{(1,1)\}
        &= \left\{(z,z):z\in\C\right\}
        \subsetneq \C\oplus\C.
    \label{eq:cpgsv17_bicf_block_separation_identical_block_diagonal_span}
\end{align}
Taking $\Delta_1=1$ and $\Delta_2=-1$ makes the trace pairing vanish for every
word and every length, contrary to
Eq.~\ref{eq:cpgsv17_bicf_block_separation_trace_separation}.

The two blocks are identical and hence related by similarity and phase. They
therefore cannot be distinct members of a minimal basis of normal tensors in
the source's sense \cite{Cirac2016MPDO_arXiv}. This example explains why the
older structure stores the finite-length witness as an independent field. It
is not a counterexample to the source proposition, and no source
counterexample is recorded.

\section{Consequence for Lemma L}

The same independence restricts the older theorem
\leanid{MPOTensor.blockwise_insert_eq_of_mpv_agree}, which assumes \leanid{MPOTensor.HorizontalCFData.biCF} for the
full family of flattened copies. Lemma~L in Appendix~C.3 of the source instead
decomposes the canonical-form tensor as
\begin{align}
    A
        &= \bigoplus_{j,q}
            \mu_{j,q}X_{j,q}A_jX_{j,q}^{-1},
    \label{eq:cpgsv17_bicf_block_separation_sector_decomposition}
\end{align}
where $j$ indexes minimal representatives and $q$ indexes their copies
\cite[lines~1835--1858]{Cirac2016MPDO_arXiv}. For each representative, the
argument groups the copies into the coefficient
\begin{align}
    c_j^{(N)}
        &= \sum_q \mu_{j,q}^{N}.
    \label{eq:cpgsv17_bicf_block_separation_grouped_power_sum}
\end{align}
The nonvanishing power-sum lemma then provides a length $N$ for which
$c_j^{(N)}\ne0$ \cite{Cirac2016MPDO_arXiv}.

This representative-grouped argument is now formalized. At a fixed length,
simultaneous separation of the representative word tuples yields
\begin{align}
    c_j^{(N)}\left(YA_j-ZA_j\right)
        &=0
        \qquad\text{for every }j.
    \label{eq:cpgsv17_bicf_block_separation_grouped_insertion_equality}
\end{align}
Choosing a length at which the bounded power sums in
Eq.~\ref{eq:cpgsv17_bicf_block_separation_grouped_power_sum} do not vanish
reduces
Eq.~\ref{eq:cpgsv17_bicf_block_separation_grouped_insertion_equality} to
\begin{align}
    YA_j
        &= ZA_j
        \qquad\text{for every }j.
    \label{eq:cpgsv17_bicf_block_separation_representative_insertion_equality}
\end{align}
A second theorem proves the required eventual simultaneous word-tuple
separation once the BNT representatives have been blocked to become
injective. Thus repeated copies and unequal nonzero copy weights no longer
obstruct the formal proof of Lemma~L.

The first-site action has also been transported through physical blocking.
Let $I=(i_0,\ldots,i_L)$ and $J=(j_0,\ldots,j_L)$ be blocked physical indices.
The induced first-site matrix is
\begin{align}
    Y^{[L+1]}_{I,J}
        &=Y_{i_0,j_0}
          \prod_{k=1}^{L}\delta_{i_k,j_k},
    \label{eq:cpgsv17_bicf_block_separation_blocked_first_site_matrix}
\end{align}
and its insertion satisfies
\begin{align}
    \left(Y^{[L+1]}A^{[L+1]}\right)^I
        &=(YA)^{i_0}A^{i_1}\cdots A^{i_L}.
    \label{eq:cpgsv17_bicf_block_separation_blocked_insertion}
\end{align}
Therefore, an all-length first-site equality before blocking implies the
corresponding blocked equality. Conversely, if the length-$L$ words of $A$
span the full matrix algebra, equality of the two blocked insertions implies
$YA=ZA$. These statements provide the transport needed to combine the source
block-injectivity argument with Lemma~L. The converse uses only block
injectivity and introduces no further hypothesis.

\section{Classification}

The scalar example in
Eq.~\ref{eq:cpgsv17_bicf_block_separation_identical_block_diagonal_span} is
not a source erratum. It shows only that per-copy blockwise injectivity,
left-canonicality, and nonzero weights do not imply simultaneous separation.
It fails the source's minimal-representative condition
\cite{Cirac2016MPDO_arXiv}.

The representative-level source theorem is now derived from the canonical-form
hypotheses. For a tensor $A$, let $S_L(A)$ denote the span of its words of
length $L$. Positive-length block injectivity propagates from $L$ to $L+1$
without a normalization assumption because
\begin{align}
    S_{L+1}(A)
        &= \operatorname{span}\{A_i\}_i S_L(A).
    \label{eq:cpgsv17_bicf_block_separation_word_span_propagation}
\end{align}
Applying the same identity at length $L$ proves the propagation result.

Physical blocking has also been defined directly for a sector decomposition
$P$. Blocking by $p$ sites preserves the representative and copy indices,
replaces $A_j$ by $A_j^{[p]}$, and replaces $\mu_{j,q}$ by $\mu_{j,q}^p$.
Consequently,
\begin{align}
    c_{P^{[p]},j}^{(N)}
        &= c_{P,j}^{(Np)}.
    \label{eq:cpgsv17_bicf_block_separation_blocked_power_sum}
\end{align}
The tensor constructed from the blocked decomposition has the same matrix
product vector (MPV) family as the physical blocking of the original tensor.

The representative-separation theorem does not require preservation of the
entire canonical-form predicate. Suppose $p>0$ and every $A_j^{[p]}$ is
injective. Pair separation for the original representatives at length $6p$
transports through the bijection between blocked and unblocked words to pair
separation for the blocked representatives at length $6$. It follows that
$P^{[p]}$ has simultaneous representative word-tuple span at every blocked
length $N$ satisfying
\begin{align}
    N
        &\geq 6(g-1)+1.
    \label{eq:cpgsv17_bicf_block_separation_eventual_blocked_span_length}
\end{align}

The first-site transport was initially proved under an explicit common
block-injectivity hypothesis. Suppose that a positive integer $L$ makes every
representative $L$-block injective, and define
\begin{align}
    p
        &= L+1.
    \label{eq:cpgsv17_bicf_block_separation_shifted_block_length}
\end{align}
Positive-length propagation makes each $A_j^{[L+1]}$ injective. Blocked
representative separation and the representative-grouped Lemma~L then give
equality of the induced insertions on every $A_j^{[L+1]}$. The original
length-$L$ span recovers equality of the unblocked insertions. The shift in
Eq.~\ref{eq:cpgsv17_bicf_block_separation_shifted_block_length} is essential:
recovering an insertion from a physical block of length $L+1$ uses the span of
its length-$L$ tail.

The common block-injectivity hypothesis now follows from the BNT
canonical-form predicate itself. Quantum Perron--Frobenius theory assigns a
finite peripheral period $m$ to an irreducible left-canonical representative
\cite{Wolf2012Quantum}. Its self-overlap along multiples of $m$ tends to $m$,
whereas the normalized BNT hypothesis makes the same subsequence tend to $1$.
Thus $m=1$, the transfer map is primitive, and the representative is normal.
The finitely many normal representatives have a common positive
block-injectivity length, obtained by taking a common multiple of their
individual lengths. The representative-indexed Lemma~L therefore follows
directly from the BNT canonical-form predicate.

The numerical Wielandt step has been sharpened to the exact common length used
in the source estimate. Let $D_j$ be the dimension of the $j$th representative
and $D$ the bond dimension of the flattened canonical-form tensor. Then
$D_j\leq D$. Quantum Wielandt gives an injective length $L_j\leq D_j^4$ for
the $j$th normal representative \cite{Sanz2010Quantum}. Positive-length
propagation therefore makes every representative injective at the common
length
\begin{align}
    L_0
        &= D^4.
    \label{eq:cpgsv17_bicf_block_separation_common_wielandt_length}
\end{align}
This exact common length has been formalized.\footnote{Formal declarations:
    \leanid{MPSTensor.SectorDecomposition.basisDim_le_totalDim},
    \leanid{MPSTensor.isNBlkInjective_pow_four_of_isNormal_leftCanonical}, and
    \leanid{MPSTensor.IsBNTCanonicalForm.basis_isNBlkInjective_totalDim_pow_four}.}

The direct-sum interpolation is also complete. Let $g$ be the number of
minimal representatives and use $L_0$ from
Eq.~\ref{eq:cpgsv17_bicf_block_separation_common_wielandt_length}. Define the
pair-span length by
\begin{align}
    S
        &= 3(L_0+1).
    \label{eq:cpgsv17_bicf_block_separation_pair_span_length}
\end{align}
The two-block argument gives full pair span at this length. Fix a block $k$,
choose an anchor block $j\ne k$, and use the pair span to realize an arbitrary
pair $(X_k,0)$. Multiplication by one identity--zero selector for each of the
remaining $g-2$ blocks produces a tuple equal to $X_k$ on block $k$ and zero
on every other block. Its length is
\begin{align}
    (g-1)S
        &=3(g-1)(D^4+1).
    \label{eq:cpgsv17_bicf_block_separation_simultaneous_interpolation_length}
\end{align}
Summing these block-supported tuples over $k$ realizes every element of
$\mathcal A_{\mathrm{prod}}$. This construction uses no separate injective
prefix. When no selectors remain, the length-zero identity is only the neutral
element of the finite product; it is neither a padding hypothesis nor a
separate theorem case.

If $g=1$, ordinary block injectivity gives the required span at length $D^4$.
If $g\geq2$, positivity of the representative dimensions and multiplicities
gives $g\leq D$. Hence
\begin{align}
    3(D^4+1)(D-1)
        &=3(D^5-D^4+D-1),
    \label{eq:cpgsv17_bicf_block_separation_bound_expansion}\\
    3(D^5-D^4+D-1)
        &\leq 3D^5,
    \label{eq:cpgsv17_bicf_block_separation_three_d_pow_five_bound}
\end{align}
because $D-1\leq D^4$.

The formal proof is recorded by the following declarations:
\begin{itemize}
    \item
    \leanid{MPSTensor.exists_mem_span_wordTuple_eq_and_eq_zero_of_pairWordTupleSpanTop}
    realizes an arbitrary value on the selected anchor pair;
    \item
    \leanid{MPSTensor.wordTupleSpanTop_mul_pred_of_forall_pairWordTupleSpanTop}
    and
    \leanid{MPSTensor.wordTupleSpanTop_mul_pred_of_forall_pairTraceSeparatingAt}
    establish direct-sum interpolation at length $(g-1)S$;
    \item
    \leanid{MPSTensor.wordTupleSpanTop_threeBlock_mul_pred_of_blocksNotGaugePhaseEquiv_c1}
    specializes the interpolation result to the source length
    $3(g-1)(L_0+1)$;
    \item
    \leanid{MPSTensor.SectorDecomposition.basisCount_le_totalDim}
    proves $g\leq D$;
    \item
    \leanid{MPSTensor.IsBNTCanonicalForm.exists_basis_wordTupleSpanTop_le_three_totalDim_pow_five}
    gives a positive simultaneous-span length bounded by $3D^5$.
\end{itemize}

The public horizontal canonical-form predicate now uses the
representative-indexed formulation. It quantifies a BNT sector decomposition
whose total bond dimension equals the original bond dimension and whose
doubled-index tensor is related to the original tensor by one global gauge.
Complete MPV equality transports first-site identities to the BNT
representatives. The representative-grouped Lemma~L gives equality of the
opposite-corner insertions on every minimal representative. Linearity extends
this equality to the full weighted direct sum, and gauge conjugation transports
it to the original matrix product operator tensor. The older per-copy
\leanid{MPOTensor.HorizontalCFData} theorems remain available only as explicitly
restricted auxiliary statements.

\section{The Appendix C.2 sector entropy argument}

The one-letter simultaneous-span assumption in the formal sector entropy
theorem is not an extra hypothesis relative to the cited source. At the start
of Case~II, the source fixes a simple tensor $K$ in block-injective canonical
form \cite[line~1628]{Cirac2016MPDO_arXiv}. Definition~2.6 of the source states
that the one-letter tuples of the minimal normal representatives span
\begin{align}
    \bigoplus_{j=1}^{g}M_{D_j}(\C).
    \label{eq:cpgsv17_bicf_block_separation_one_letter_bicf_span}
\end{align}
The inverse tensor used immediately afterward, and hence the one-site
projections in the proof of Proposition~4.13, depend on this hypothesis
\cite{Cirac2016MPDO_arXiv}. The formal assumption
\leanid{MPSTensor.WordTupleSpanTop S.basis 1} is therefore the direct algebraic form of
the source's block-injective canonical-form hypothesis.

Starting instead from an unblocked canonical form and removing the one-letter
span by choosing a finite blocking length would be strictly stronger. Although
the sector area law (SAL) passes from a tensor to every positive physical
blocking, the periodic block identity does not provide the converse. An
$L$-blocked chain records only original chain lengths of the form $NL$ and
only cuts at multiples of $L$. In particular, SAL for $K^{[L]}$ does not by
itself recover
\begin{align}
    I_m(K)
        &= I_{m+1}(K),
    \label{eq:cpgsv17_bicf_block_separation_unblocked_entropy_equality}
\end{align}
for arbitrary chain lengths and cuts. The proposed reverse transport in
issue~\#4119 is therefore unavailable without a new argument absent from
Appendix~C.2, and the printed proposition does not require such an argument
\cite{Cirac2016MPDO_arXiv}.

The sector entropy theorem should consequently retain the simultaneous
one-letter span and identify it with the source block-injective canonical-form
hypothesis. The finite-blocking theorem shows that an arbitrary canonical-form
tensor can be placed in block-injective canonical form when the surrounding
result is itself stated only up to physical blocking. It does not convert the
Appendix~C.2 conclusion into an unblocked SAL statement.

The representative argument closes the multiplicity and block-separation gap
in Lemma~L. It does not by itself complete Proposition~4.13 of the source:
constructing the letterwise vertical decomposition from the resulting reducing
projections, and then proving positivity and isometric normalization of the
sectors, remain separate parts of that proposition
\cite{Cirac2016MPDO_arXiv}.

The gap is resolved. The formal theorem now matches the source's hypotheses and
conclusion through a different proof of simultaneous interpolation. BNT
canonical form supplies normality, left-canonicality, normalized self-overlap,
and pairwise gauge-phase inequivalence; no finite separation witness is added
as a hypothesis. The source-level Lemma~L requires only the existence of a
suitable finite length, while the formal theorem also proves the stated
$3D^5$ bound.

\bibliographystyle{alpha}
\bibliography{references}

\end{document}
